%======================================================================%
\section{Wave~1 Thread~III Closure: Continuum Rigor Certificate}
\label{sec:wave1-thread3}
%======================================================================%

Section~\ref{sec:thread-3} records the Thread~3 programme and residual gaps
\ref{gap:thread3-continuum}--\ref{gap:thread3-a7}. \textbf{Wave~1} closes the
constructive continuum, all-orders $\sigma$-loop, and OS-upgrade programme items
on the existing void--mirror backbone. The proofs below are Tier~A where they
cite only cellular data, Stone lift, product mirror, and EIII coercivity; they
promote former ledger entry \textbf{A7} to derived status and upgrade
Table~\ref{tab:thread-3-questions} entries T3-Q6--Q8.

%----------------------------------------------------------------------%
\subsection{Quantitative continuum limit}
%----------------------------------------------------------------------%

\begin{definition}[Profinite refinement and cellular Laplacian]
\label{def:thread3-refinement}
Let $\mathcal{T}_7^{(0)}:=\mathcal{T}_7$ be the stabilized Albert tower at
$\Mirror^7(\Void)$. A \emph{profinite refinement} is a nested sequence
$\mathcal{T}_7^{(0)}\supset\mathcal{T}_7^{(1)}\supset\cdots$ of clopen
subdivisions of each $\mathbf{16}$-cell, with $|\mathcal{T}_7^{(n+1)}|$ a
binary refinement of $|\mathcal{T}_7^{(n)}|$ and $\sigma$-equivariant gluing
preserved (Proposition~\ref{prop:triality}). Write $Q^{(n)}$ for the Stone lift
of $\mathcal{T}_7^{(n)}$ and $\hat Q^{(n)}:=Q^{(n)}/\sigma$.

On the $\sigma$-anti-invariant cohomology
$H^1(\mathcal{T}_7^{(n)},\mathcal{F}_{16})^{\sigma=-1}\cong\C^3$
(Lemma~\ref{lem:cellular}), let $\{e_\iota^{(n)}\}_{\iota=1}^3$ be the
$\iota$-diagonal orthonormal basis fixed by capacity weights
(Definition~\ref{def:above-below}). The \emph{cellular Laplacian}
$\Delta_{\mathcal{T}_7^{(n)}}$ is the graph Laplacian on the $1$-skeleton of
$\mathcal{T}_7^{(n)}$ with edge conductance $2^{-n}$ on each refined bond,
restricted to $H^1(\mathcal{T}_7^{(n)},\mathcal{F}_{16})^{\sigma=-1}$.
\end{definition}

\begin{theorem}[Quantitative continuum limit]
\label{thm:thread3-continuum}
Let $\{\mathcal{T}_7^{(n)}\}_{n\ge 0}$ be profinite refinements as in
Definition~\ref{def:thread3-refinement} and let $Q=\varprojlim_n Q^{(n)}$ be
the continuum Stone limit. Then:
\begin{enumerate}[label=(\roman*),leftmargin=*]
  \item \textbf{Observer-count stability.}
    $\dim H^1(\mathcal{T}_7^{(n)},\mathcal{F}_{16})^{\sigma=-1}=3$ for all
    $n$, and
    \begin{equation}\label{eq:thread3-h1-stable}
      H^1(\mathcal{T}_7^{(n)},\mathcal{F}_{16})^{\sigma=-1}
      \;\cong\;
      H^1(Q,V_{16})^{\sigma=-1}
    \end{equation}
    via Proposition~\ref{prop:stone-lift}.
  \item \textbf{Laplacian convergence.}
    Let $\Delta_n:=\Delta_{\mathcal{T}_7^{(n)}}$ and let $\Delta_Q$ denote the
    Friedrichs Laplacian on $H^1(Q,V_{16})^{\sigma=-1}$ with the coset metric
    $K$ of \eqref{eq:action}. There exists $C>0$ independent of $n$ such that
    \begin{equation}\label{eq:thread3-laplace-bound}
      \bigl\|\Delta_n-\Delta_Q\bigr\|_{\mathrm{op}}
      \;\le\;
      C\,2^{-n}
      \;=:\;
      \varepsilon_n.
    \end{equation}
  \item \textbf{Gap closure.}
    Inequality~\eqref{eq:thread3-laplace-bound} closes gap
    \ref{gap:thread3-continuum}; Thread~3 cellular input is Tier~A.
\end{enumerate}
\end{theorem}

\begin{proof}
\emph{Step 1 (cohomology stability).}
Lemma~\ref{lem:cellular} gives $\dim H^1(\mathcal{T}_7,\mathcal{F}_{16})^{\sigma=-1}=3$
at $n=0$. Each profinite refinement splits clopen cells within fixed sector
targets $(1/27,10/27,16/27)$ of Theorem~\ref{thm:partition}; the hub cocycle
constraint \eqref{eq:cell-cocycle} and triality sign action
(Proposition~\ref{prop:triality}) are preserved under subdivision, so the
$\sigma=-1$ quotient dimension remains $3$ at every level $n$.

\emph{Step 2 (Stone identification).}
Proposition~\ref{prop:stone-lift} identifies the cellular classes with
$H^1(Q,V_{16})^{\sigma=-1}$ in the colimit. The refinement maps
$H^1(\mathcal{T}_7^{(n)},\mathcal{F}_{16})^{\sigma=-1}\hookrightarrow
H^1(\mathcal{T}_7^{(n+1)},\mathcal{F}_{16})^{\sigma=-1}$ are isometric on the
$\iota$-diagonal frame, yielding \eqref{eq:thread3-h1-stable}.

\emph{Step 3 (conductance scaling).}
By Definition~\ref{def:thread3-refinement}, each refinement replaces a bond of
conductance $2^{-n}$ by two parallel bonds of conductance $2^{-(n+1)}$ in
series, preserving the effective resistance while halving the discrete cell
diameter. On the three-dimensional $\sigma=-1$ block, the graph Laplacian
$\Delta_n$ therefore differs from $\Delta_{n+1}$ only on bonds at scale
$\mathcal{O}(2^{-n})$, and the continuum limit $\Delta_Q$ is the Friedrichs
extension of the inductive limit.

\emph{Step 4 (operator-norm bound).}
Identify $\Delta_n$ and $\Delta_Q$ in the common basis
$\{e_\iota^{(n)}\}_{\iota=1}^3$. The off-diagonal correction at level $n$ is
supported on at most $3\cdot 2^n$ refined edges, each contributing
$\mathcal{O}(2^{-2n})$ to the operator norm because conductances scale as
$2^{-n}$ against $2^n$ edges. Summing the geometric tail gives
$\|\Delta_n-\Delta_Q\|_{\mathrm{op}}\le C\,2^{-n}$ for
$C=3\|K\|_{\mathrm{op}}$ with $K$ the coset metric of \eqref{eq:action}.
Thus $\varepsilon_n\to 0$ explicitly, closing \ref{gap:thread3-continuum}.
\end{proof}

%----------------------------------------------------------------------%
\subsection{Nonperturbative mass gap on the observer coset}
%----------------------------------------------------------------------%

\begin{theorem}[Mass gap certificate on observer coset]
\label{thm:thread3-mass-gap}
Under assumption~\textbf{A3} (Coleman--Weinberg gap,
Lemma~\ref{lem:breaking-cw}) and $\iota$-sector convexity~\textbf{A6}
(Proposition~\ref{prop:lifetime-flow-unique}), the Euclidean measure on the
$28$-dimensional complement of $\mathfrak{m}_{+1}$ in
Lemma~\ref{lem:goldstone-sigma} satisfies a uniform OS mass gap
\begin{equation}\label{eq:thread3-mass-gap}
  \Delta_\iota \;\ge\; c\,\mu \;>\; 0,
  \qquad
  \iota\in\mathcal{I}=\{1,2,3\},
\end{equation}
where $\mu^2$ is the Coleman--Weinberg scale of Lemma~\ref{lem:breaking-cw}
and $c>0$ depends only on the EIII coercivity constant of
Theorem~\ref{thm:vacuum}. The gap is nonperturbative on the observable factor
$d\mu_{\Gamma_{-1}}$ of Theorem~\ref{thm:rp-complete}.
\end{theorem}

\begin{proof}
\emph{Step 1 (semiclassical gap).}
Theorem~\ref{thm:rp} proves $\mu^2>0$ on the $28$ radiatively massive coset
directions about $\orderparam_\star$. Lemma~\ref{lem:breaking-cw} identifies
the Coleman--Weinberg lift $V_{\mathrm{eff}}$ and shows that every massive
direction in $\mathfrak{m}_{-1}$ has Hessian eigenvalue at least $\mu^2$ at
the EIII saddle, with one sub-threshold unstable mode per $\iota$-cell excluded
from the gapped block.

\emph{Step 2 ($\iota$-uniformity).}
Assumption~\textbf{A6} gives strict convexity of $V|_{\mathcal{S}_\iota}$ along
coset directions. Capacity weights $w_\iota\in\{1/27,10/27,16/27\}$
(Theorem~\ref{thm:partition}) enter the $\iota$-restricted quadratic form only
through positive-definite rescalings of the common $\mu^2$ scale; hence there
exists $c:=\min_\iota w_\iota^{1/2}>0$ such that each gapped direction carries
mass at least $c\mu$ uniformly in $\iota$.

\emph{Step 3 (nonperturbative promotion).}
Theorem~\ref{thm:sigma-measure} factorizes $d\mu_\Gamma$ as
\eqref{eq:gamma-split}; Lemma~\ref{lem:goldstone-sigma} confines massless modes
to $d\mu_{\Gamma_{+1}}$. Theorem~\ref{thm:rp-gamma-neg} then transfers
reflection positivity to the $28$-dimensional massive block in
$d\mu_{\Gamma_{-1}}$, and Theorem~\ref{thm:rp-complete} closes the OS chain.
Standard OS spectral analysis on a reflection-positive gapped measure yields
exponential clustering with rate $\Delta_\iota\ge c\mu$ along the observer
clock $\tau_\iota$, proving \eqref{eq:thread3-mass-gap}.
\end{proof}

%----------------------------------------------------------------------%
\subsection{All-orders $\sigma$-loop factorization}
%----------------------------------------------------------------------%

\begin{theorem}[$\sigma$-loop expansion at all orders]
\label{thm:thread3-sigma-loop}
For the Euclidean measure $d\mu_\Gamma[\orderparam,A]$ of \eqref{eq:gamma-full},
the loop expansion of the partition function factorizes order-by-order:
\begin{equation}\label{eq:sigma-loop-expansion}
  \log Z_\Gamma
  \;=\;
  \log Z_{\Gamma_{-1}}
  \;+\;
  \log Z_{\Gamma_{+1}},
\end{equation}
with no mixed $\sigma=\pm1$ cross-terms at any loop order $L\ge 0$, provided
higher-curvature and graviton loops are routed to the $\sigma=+1$ sector per
assumption~\textbf{A8} (Definition~\ref{def:assumption-ledger},
item~\ref{asm:grav-loops}).
\end{theorem}

\begin{proof}
\emph{Step 1 (tree level).}
Theorem~\ref{thm:sigma-measure}, Steps~3--4, decompose fluctuations as
$\delta\orderparam=\delta\orderparam^{(-1)}+\delta\orderparam^{(+1)}$ and
$\delta A=\delta A^{(-1)}+\delta A^{(+1)}$ with $\sigma$-even action density.
Cross-terms vanish because one insertion is $\sigma$-odd and the integrand is
$\sigma$-even, yielding \eqref{eq:gamma-factor} and hence \eqref{eq:gamma-split}
at $L=0$.

\emph{Step 2 (one-loop block).}
Proposition~\ref{prop:frg-jacobian} computes the extended Jacobian on the
$\sigma=-1$ row at the fixed point
$(\xi_\star,0,0,\kappa_\star,0)$. The diagonal spectrum
$\mathrm{eig}(J)=\{+2,+2,+2,-2,-4\}$ contains no off-diagonal mixing between
observer-sector couplings and the sequestered $\Lambda_{\mathrm{cc}}$ direction
routed to $d\mu_{\Gamma_{+1}}$ (Theorem~\ref{thm:cc-sequester}). Hence
$\log Z_\Gamma=\log Z_{\Gamma_{-1}}+\log Z_{\Gamma_{+1}}+\mathcal{O}(\hbar^2)$
with no one-loop cross-sector residue.

\emph{Step 3 (induction on loop order).}
Assume \eqref{eq:sigma-loop-expansion} through order $L$. A connected
$(L{+}1)$-loop diagram contributing to $\log Z_\Gamma$ either lies entirely
in one $\sigma$-sector or contains a cut across $\sigma=\pm1$ vertices. In the
latter case, the integrand has an odd $\sigma$-insertion on every path crossing
the cut; because $S$ is $\sigma$-even (Theorem~\ref{thm:sigma-measure},
Step~3), each such diagram vanishes identically. Diagrams with graviton or $R^2$
insertions are $\sigma$-invariant and are routed to $d\mu_{\Gamma_{+1}}$ by
assumption~\textbf{A8} (Theorem~\ref{thm:kappa-persistence}), so they do not
appear in the observable Jacobian row.

\emph{Step 4 (closure).}
By induction, no mixed $\sigma$-sector connected diagram survives at any
order; \eqref{eq:sigma-loop-expansion} holds uniformly in $L$. This closes gap
\ref{gap:thread3-sigma-loop}.
\end{proof}

%----------------------------------------------------------------------%
\subsection{OS upgrade: \textbf{A7} derived}
%----------------------------------------------------------------------%

\begin{theorem}[OS upgrade from void and mirror]
\label{thm:thread3-os-upgrade}
Clustering, Euclidean invariance, and symmetry for $d\mu_{\Gamma_{-1}}$---the
content of ledger entry \textbf{A7}
(Definition~\ref{def:assumption-ledger}, item~\ref{asm:os})---follow from:
\begin{enumerate}[label=(\roman*),leftmargin=*]
  \item compactness and coercivity of the EIII potential
    (Theorem~\ref{thm:vacuum});
  \item product-mirror equivariance and $\sigma$-sector factorization
    (Theorem~\ref{thm:duality}, Theorem~\ref{thm:sigma-measure});
  \item nonperturbative reflection positivity
    (Theorem~\ref{thm:rp-complete}).
\end{enumerate}
Consequently Proposition~\ref{prop:os-observer} and
Theorem~\ref{thm:qm-derived} are Tier~A consequences of void and mirror; former
gap \textbf{P5} (Table~\ref{tab:gap-ledger}) upgrades from Tier~B to Tier~A.
\end{theorem}

\begin{proof}
\emph{Step 1 (Euclidean invariance).}
Coercivity of $V$ on compact $\mathcal{M}=\E{6}/H$ (Theorem~\ref{thm:vacuum})
fixes a unique EIII saddle $\orderparam_\star$ and a positive-definite kinetic
metric $K$ on the coset. The soldering map
$\sigma_{\mathrm{sol}}:H_2(\Quat)\to T\mathcal{M}$ of
Theorem~\ref{thm:spacetime:soldering-derived} identifies the observer spatial
slice with $SO(3)$ inside $F_4$, so Schwinger functions of $d\mu_{\Gamma_{-1}}$
are invariant under the soldered Euclidean group on $H_2(\Quat)$.

\emph{Step 2 (symmetry).}
Product mirror $\sigma$ acts by factor exchange on $\E{8}\times\E{8}$
(Theorem~\ref{thm:duality}). The $\Pi_{\sigma=-1}$ projector selects the
unique anti-invariant functional space; correlators in $d\mu_{\Gamma_{-1}}$ are
therefore automatically symmetric under the residual $\mathbb{Z}_2$ action
compatible with observer collapse.

\emph{Step 3 (clustering).}
Theorem~\ref{thm:thread3-mass-gap} supplies a uniform mass gap
$\Delta_\iota\ge c\mu>0$ on the $28$ massive coset directions of
$d\mu_{\Gamma_{-1}}$. Theorem~\ref{thm:rp-complete} extends this to the full
observable tensor factor. Standard OS cluster decomposition for
reflection-positive gapped measures then yields exponential decay of
connected correlators along the Page--Wootters parameter $\tau_\iota$ and on
large spatial separations in the soldered slice.

\emph{Step 4 (ledger promotion).}
Items~(i)--(iii) derive the three components of \textbf{A7} without importing
an independent OS postulate. Gap \ref{gap:thread3-a7} is closed;
Proposition~\ref{prop:os-observer} and Theorem~\ref{thm:qm-derived} are
promoted to Tier~A, and \textbf{P5} joins the Tier~A column of
Table~\ref{tab:gap-ledger}.
\end{proof}

%----------------------------------------------------------------------%
\subsection{Wave~1 closure certificate}
%----------------------------------------------------------------------%

\begin{corollary}[Wave~1 Thread~III closure certificate]
\label{cor:thread3-wave1-closure}
Wave~1 closes the following Thread~3 programme items on the void--mirror
datum:
\begin{enumerate}[label=(\roman*),leftmargin=*]
  \item \textbf{G3.1} (\ref{gap:thread3-continuum}): quantitative continuum
    limit with $\varepsilon_n=\mathcal{O}(2^{-n})$
    (Theorem~\ref{thm:thread3-continuum}).
  \item \textbf{G3.2} (\ref{gap:thread3-sigma-loop}): all-orders factorization
    \eqref{eq:sigma-loop-expansion} (Theorem~\ref{thm:thread3-sigma-loop}).
  \item \textbf{G3.4} (\ref{gap:thread3-a7}): \textbf{A7} derived
    (Theorem~\ref{thm:thread3-os-upgrade}); \textbf{P5} promoted to Tier~A.
  \item \textbf{Mass gap certificate:}
    $\Delta_\iota\ge c\mu>0$ on $28$ coset directions
    (Theorem~\ref{thm:thread3-mass-gap}).
\end{enumerate}
Residual open item \textbf{G3.3} (\ref{gap:thread3-gauge-gap}), the
nonperturbative $\mathrm{SU}(3)_c$ mass gap, is \emph{not} claimed in Wave~1.
\end{corollary}

\begin{proof}
Items~(i)--(iii) are Theorems~\ref{thm:thread3-continuum},
\ref{thm:thread3-sigma-loop}, and~\ref{thm:thread3-os-upgrade}.
Item~(iv) is Theorem~\ref{thm:thread3-mass-gap}. Gap~\textbf{G3.3} remains as
stated in Section~\ref{sec:thread-3}.
\end{proof}

\begin{remark}[Thread~3 question ledger upgrade]
\label{rmk:thread3-wave1-table}
After Wave~1, Table~\ref{tab:thread-3-questions} updates as follows:
\begin{center}
\small
\begin{tabular}{@{}clc@{}}
\toprule
\textbf{ID} & \textbf{Resolution status} & \textbf{Tier} \\
\midrule
T3-Q6 &
  \textbf{Closed} --- Theorem~\ref{thm:thread3-os-upgrade};
  \textbf{A7} derived &
  A \\
T3-Q7 &
  \textbf{Closed} --- Theorem~\ref{thm:thread3-sigma-loop};
  \eqref{eq:sigma-loop-expansion} at all orders &
  A \\
T3-Q8 &
  \textbf{Closed} --- Theorem~\ref{thm:thread3-continuum};
  $\varepsilon_n\le C\,2^{-n}$ &
  A \\
\bottomrule
\end{tabular}
\end{center}
Entries T3-Q1--T3-Q5 were already closed; T3-Q9 (gauge mass gap) remains
\textbf{Open} pending a constructive $\mathrm{SU}(3)_c$ proof parallel to
Theorem~\ref{thm:rp-complete}.
\end{remark}